\subsection{Extension of a hypocontinuous bilinear mapping}

PROPOSITION 7. --- Let E, F, G be three locally convex spaces, G being assumed Hausdorff; let \(E_{0}\) (resp. \(F_{0}\) ) be a dense vector subspace of E (resp. F). Let u be a separately continuous bilinear mapping from \(E \times F\) into G.

\begin{enumerate}
\def\labelenumi{\arabic{enumi})}
\item
  If \(u(\mathrm{E}_0 \times \mathrm{F}_0) = \{0\}\) , then \(u = 0\) .
\item
  Let \(\mathfrak{S}_0\) be a family of bounded subsets of \(\mathbf{E}_0\) ; if the restriction of \(u\) to \(\mathbf{E}_0 \times \mathbf{F}_0\) is \(\mathfrak{S}_0\) -hypocontinuous then so is \(u\) .
\item
  By hypothesis, for all \(x \in \mathrm{E}_0\) , the continuous linear mapping \(u(x, .)\) is null on \(\mathrm{F}_0\) , hence on \(\mathrm{F}\) : therefore for all \(y \in \mathrm{F}\) , the continuous linear mapping \(u(., y)\) is null on \(\mathrm{E}_0\) , hence on \(\mathrm{E}\) . This proves that \(u = 0\) .
\item
  For every closed neighbourhood W of 0 in G and for every set \(\mathbf{M} \in \mathfrak{S}_0\) , there exists, by hypothesis, a neighbourhood V of 0 in \(\mathrm{F}_0\) such that \(u(\mathbf{M} \times \mathbf{V}) \subset \mathbf{W}\) . But \(\overline{\mathbf{V}}\) is a neighbourhood of 0 in F; for every \(x \in \mathbf{M}\) , the relation \(u(\{x\} \times \mathbf{V}) \subset \mathbf{W}\) implies that \(u(\{x\} \times \overline{\mathbf{V}}) \subset \mathbf{W}\) , since \(u(x, .)\) is continuous and W is closed; therefore \(u(\mathbf{M} \times \overline{\mathbf{V}}) \subset \mathbf{W}\) , which proves that \(u\) is \(\mathfrak{S}_0\) -hypocontinuous.
\end{enumerate}

PROPOSITION 8. --- Let E, F, G be three locally convex spaces; assume that G is Hausdorff and quasi-complete. Let \(E_{0}\) (resp. \(F_{0}\) ) be a dense vector subspace of E (resp. F), \(S_{0}\) (resp. \(T_{0}\) ) a family of bounded subsets of \(E_{0}\) (resp. \(F_{0}\) ) such that every point of E (resp. F) is in the closure of an element of \(S_{0}\) (resp. \(T_{0}\) ). Then every \((\mathfrak{S}_{0}, \mathfrak{T}_{0})\) -hypocontinuous bilinear mapping u from \(E_{0} \times F_{0}\) into G extends uniquely to a separately continuous bilinear mapping \(\overline{u}\) from E × F into G and \(\overline{u}\) is \((\mathfrak{S}_{0}, \mathfrak{T}_{0})\) -hypocontinuous.

The uniqueness and hypocontinuity of \(\overline{u}\) follows from prop. 7; it remains to prove the existence of \(\overline{u}\) . For every \(y' \in F_{0}\) , the continuous linear mapping \(x' \mapsto u(x', y')\) from \(E_{0}\) into G extends uniquely to a continuous linear mapping \(x \mapsto u_{1}(x, y')\) from E into G (III, p.~8, prop. 10). It follows immediately that for every \(x \in E\) , the mapping \(y' \mapsto u_{1}(x, y')\) from \(F_{0}\) into G is linear; and we shall show that it is continuous. By hypothesis, there exists \(M \in S_{0}\) , such that \(x \in M\) . For every closed neighbourhood W of 0 in G, there exists, by hypothesis, a neighbourhood V of 0 in \(F_{0}\) such that \(u(M \times V) \subset W\) ; since \(x \mapsto u_{1}(x, y')\) is continuous, we deduce that \(u_{1}(\overline{\mathbf{M}} \times \mathbf{V}) \subset \mathbf{W}\) , and in particular \(u_{1}(x, y') \in \mathbf{W}\) for all \(y' \in V\) . This establishes our assertion. By virtue of prop. 7, the bilinear map \(u_{1}\) from \(E \times F_{0}\) into G is \((\mathfrak{S}_{0}, \mathfrak{T}_{0})\) -hypocontinuous. We end the proof by interchanging the roles of E and F in the first part of the proof, applied to \(u_{1}\) .
% semantic-fragments: legacy-input-seam
\relax{}
